\documentclass[11pt,a4paper]{book}

\usepackage[utf8]{inputenc}

\usepackage[T1]{fontenc}

\usepackage[french]{babel}

\usepackage{amsmath,amssymb,amsthm}

\usepackage{tikz}

\usepackage{tikz-3dplot}

\usepackage{pgfplots}

\pgfplotsset{compat=1.18}

\usetikzlibrary{positioning,shapes.geometric,arrows.meta,decorations.pathreplacing}

\usepackage{tcolorbox}

\usepackage{hyperref}

\usepackage{xcolor}

\usepackage{mathrsfs}

\usepackage{algorithm}

\usepackage{algpseudocode}

\begin{document}

\chapter{Hypotheses emises Ordonnance de renvoie: }

\section{hypothese 1 : Suicide}

\textbf{Hypothese refutée} car il n'a pas pu s'infliger lui-meme les sevices corporel qu'il a subit avant sa mort.

\section{hypothese 2: Le coupable n'est pas parmi les personnes qui sont inculpées}

\textbf{Hypothese refutée} car les dernieres personnes presentes avec lui lors de sa mort etaient les personnes inculpees. Il n'ya donc pas de personnage tier dans cette affaire.

\section{hypothese 3: Le coupable est parmi les personnes qui sont inculpées}

\textbf{Hypothese coherente :} Les donnees de localisation croisees aux differents temoignages prouvent que certaines des personnes inculpees etaient les dernieres presentes avec la victime avant sa mort.

\end{document}